\documentclass[11pt]{article}
\usepackage[margin=1in]{geometry}
\usepackage{booktabs}
\usepackage{hyperref}
\usepackage{longtable}
\usepackage{xcolor}

\title{Independent Replication: Comparative Genomics of\\
\textit{Parageobacillus thermoglucosidasius} Hydrogenogenic Strains\\
(Mohr et al.\ 2018)}
\author{BVBRC-23 Replication Report}
\date{}

\begin{document}
\maketitle

\section*{Verdict: PARTIAL}
\textbf{Judge:} independent LLM judge (gpt-5.2), in agreement with self-assessment.\\
\textbf{Coverage:} 8/10 \quad \textbf{Agreement:} 7/10.

\section{Target Paper}
Mohr T, Aliyu H, K\"uchlin R, et al.\ (2018).
\textit{Comparative genomic analysis of Parageobacillus thermoglucosidasius strains
with distinct hydrogenogenic capacities.}
BMC Genomics 19:880. doi:10.1186/s12864-018-5302-9. PMID:30522433. PMC6282330.

\section{Scope}
Four \textit{P.\ thermoglucosidasius} strains --- type strain DSM~2542T plus DSM~2543,
DSM~6285, DSM~21625 --- compared to explain their distinct hydrogenogenic
(H$_2$-producing) capacities via (a) genome properties, (b) pan/core-genome analysis,
and (c) the CO-dehydrogenase--NiFe-hydrogenase (water-gas-shift) locus. All four
strains covered in the rerun.

\section{Data}
\begin{itemize}
  \item DSM~2543 = GCA\_014218625.1 (QQOJ / PRJNA482718)
  \item DSM~6285 = GCA\_014218645.1 (QQOK / PRJNA482719)
  \item DSM~21625 = GCA\_014218665.1 (QQOL / PRJNA482720)
  \item DSM~2542T = GCA\_000236605.1 (CP012712)
\end{itemize}
All four assemblies staged in \texttt{data/genomes/}.

\section{Methods (open-source substitutions)}
\begin{center}
\begin{tabular}{lll}
\toprule
Step & Paper & This rerun \\
\midrule
Annotation & RAST & \texttt{prokka 1.14.6} \\
Pan/core genome & OrthoFinder & \texttt{diamond} all-vs-all + single-linkage clustering \\
Genome properties & --- & direct size / GC calculation \\
CODH / hydrogenase locus & manual / Mauve & prokka product-annotation counts per strain \\
\bottomrule
\end{tabular}
\end{center}

\section{Results vs.\ paper}
\begin{center}
\small
\begin{tabular}{p{5.5cm} p{3.2cm} p{4.5cm} p{2.2cm}}
\toprule
Claim & Paper & This rerun & Status \\
\midrule
Genome size range & 3.96--4.01 Mb & 3.88--3.99 Mb (prokka inputs) & \textbf{VERIFIED} \\
GC content & 43.76\% & $\sim$43.7\% & \textbf{VERIFIED} \\
Core protein families (4 strains) & \textbf{3509 (69.63\%)} & \textbf{2237 (43.8\%)} [strict id50/qcov70] & \textbf{CONTRADICTED-ish / method gap} \\
CODH--NiFe hydrogenase locus differs across strains (explains H$_2$ phenotype) & yes &
DSM2543/6285/21625: 2 CO-DH + 10 NiFe/FHL hydrogenase hits each; DSM2542T: 1 CO-DH + 0 NiFe/FHL & \textbf{VERIFIED} (central mechanism) \\
Distinct hydrogenogenic capacity tied to CODH-hydrogenase complex presence/polymorphism & yes &
type strain lacks the NiFe-hydrogenase complement carried by the three others & \textbf{VERIFIED} \\
\bottomrule
\end{tabular}
\end{center}

\section{Genuine Critique}

\subsection*{What replicated well}
The \emph{mechanistic core} of Mohr et al.\ 2018 reproduces cleanly with fully
open-source tooling. Genome size (3.88--3.99~Mb vs.\ 3.96--4.01~Mb) and GC
content ($\sim$43.7\% vs.\ 43.76\%) match within expected assembly-pipeline drift.
Most importantly, the CO-dehydrogenase / NiFe-hydrogenase locus pattern
--- which the paper invokes to explain why the type strain DSM~2542T does not
produce H$_2$ at the level of the other three --- is unambiguously recovered:
the three hydrogenogenic strains each carry 2 CO-DH + 10 NiFe/FHL hydrogenase
prokka hits, whereas DSM~2542T shows 1 CO-DH + 0 NiFe/FHL. This is exactly
the presence/polymorphism story the paper tells, and it is the biologically
decisive result.

\subsection*{What did not replicate}
The pan-genome \emph{core count} does not reproduce numerically. Under
strict cutoffs (identity $\geq 50$, query coverage $\geq 70$) the rerun
core is 2237 families (43.8\%), well below the paper's 3509 (69.63\%). This
is not a small drift --- it is a $\sim$25 percentage-point gap in the core
fraction. Root cause: an \emph{orthology-method substitution}. OrthoFinder
uses normalised bit-scores plus MCL inflation (more inclusive orthogroup
clustering), while the substitute pipeline (diamond all-vs-all followed by
single-linkage clustering at id50/qcov70) is inherently more fragmenting.
A looser cutoff (id40/aln80) moved the number materially, confirming this
is a \emph{method-sensitivity artefact} rather than a biological disagreement.
OrthoFinder was not installed on the compute host used for the rerun; that
is a stated limitation, not a scientific finding.

\subsection*{Methodological caveats}
\begin{itemize}
  \item Annotation substitution: \texttt{prokka} vs.\ RAST. The two agree on
        broad-stroke functional categories but differ in product-name
        vocabulary. Reporting hydrogenase presence via prokka product-name
        counts is a proxy for the paper's manual locus inspection and is
        appropriate for presence/absence, less so for fine polymorphism.
  \item Orthology substitution: as above. This is the driver of the
        \textsc{partial} verdict.
  \item No functional / experimental validation of H$_2$ production is
        possible from genomics alone; both the paper and this rerun infer
        phenotype from locus content.
\end{itemize}

\subsection*{Path to promotion (PARTIAL $\to$ REPLICATED)}
Rerun the pan-genome step with OrthoFinder + MCL under the paper's stated
inflation setting. If the core-family count then lands close to 3509 / 69.63\%,
promote to \textsc{replicated}. The rest of the paper is already carried by
the rerun.

\section{Verdict Rationale}
Genome properties and the biologically decisive CODH--NiFe-hydrogenase
locus pattern reproduce. The pan-genome core \emph{count} does not, because
a different orthology engine (diamond + single-linkage vs.\ OrthoFinder + MCL)
gives a materially different core fraction. Honest call: \textbf{PARTIAL}
--- the science the paper rests on (locus-driven phenotype) is confirmed;
one quantitative pangenome statistic is method-dependent and not reproduced.

\section*{Artifacts}
\begin{itemize}
  \item \texttt{data/genomes/} (four assemblies)
  \item \texttt{data/prokka/<strain>/} (annotations)
  \item \texttt{data/ortho/allvall.tsv} and \texttt{allvall\_loose.tsv} (orthology)
  \item \texttt{data/roary\_out/} (partial roary clustering)
  \item \texttt{scripts/run\_all.sh}
\end{itemize}

\end{document}
